% =====================================================================
% Paper 06 — Pillar-6 electroweak-bridge programme note with
%            first ΔF<0 broken-energy data point at N=16
%            (Math292 4-gauge framework: G1 of 4 gauges only)
% Status: [DRAFT] [NEEDS_UPDATE] — Wave 4 generations cluster
% AUDIT-FLAG: AUDIT-2026-05-02-Wave7-Aux-Epoch-Overclaim
%   (Math314-AddC extension covering Wave 4/5 papers Paper-06..08).
% Rev 4 (2026-05-06):
%   (a) Introduction toned down: removed "three generations emerge",
%       "Yukawa hierarchy", "Higgs doublet structure cleanly". Recast
%       as exploration of an embedding programme; the broken phase is
%       a candidate Higgs-doublet support, not a verified embedding.
%   (b) Theorem H2 corrected: replaced misuse of "A-valid" with
%       "G1 only / broken-energy entry gauge". Math292 4-gauge
%       acceptance is f>0 ∧ ΔF<0 ∧ λ_min^transverse≥0 ∧ extraction_OK.
%       The current data point fulfils G2 (ΔF<0) at N=16 only;
%       G1, G3, G4 are pending re-certification per Math310-AddA.
%   (c) N=32 ΔF=-1021.7 reframed as plateau / pending verdict, NOT
%       confirmed continuum-limit endpoint. Continuum extrapolation
%       requires a second sub-N data point and a verified
%       quadratic-fit residual.
% Compile: bash _shared/build-paper.sh Paper-06-Pillar6-Generations/Paper-06
%          (or .\_shared\build-paper.ps1 ...)
% Canonical archive: Math234, Math235, Math236, Math292, Math293,
%                    Math294, Math294-AddA, Math290, Math310-AddA,
%                    Math313, Math314-AddC
% =====================================================================
% --- TECT canonical preamble (see _shared/tect-preamble.tex) ----------
% =====================================================================
% TECT canonical LaTeX preamble (revtex4-2 / single-column / theorem-safe)
% =====================================================================
% Maintainer: Jusang Lee (jtkor@outlook.com)
% Binding from: 2026-05-06
% Purpose: a single source-of-truth preamble for every TECT paper
% (Paper-00 .. Paper-10 + Auxiliary notes). Designed to (a) eliminate
% font-corruption on pdflatex (T1 fontenc + lmodern), (b) make
% theorem-heavy mathematical-physics content render cleanly in a
% single-column layout (PRD class option, NOT PRL two-column), (c)
% provide consistent theorem numbering across all TECT papers via
% amsthm with a shared counter set, (d) make hyperref + cleveref
% cross-references resolve cleanly with the standard two-pass
% pdflatex compile.
%
% Usage in each Paper-NN.tex:
%   % =====================================================================
% TECT canonical LaTeX preamble (revtex4-2 / single-column / theorem-safe)
% =====================================================================
% Maintainer: Jusang Lee (jtkor@outlook.com)
% Binding from: 2026-05-06
% Purpose: a single source-of-truth preamble for every TECT paper
% (Paper-00 .. Paper-10 + Auxiliary notes). Designed to (a) eliminate
% font-corruption on pdflatex (T1 fontenc + lmodern), (b) make
% theorem-heavy mathematical-physics content render cleanly in a
% single-column layout (PRD class option, NOT PRL two-column), (c)
% provide consistent theorem numbering across all TECT papers via
% amsthm with a shared counter set, (d) make hyperref + cleveref
% cross-references resolve cleanly with the standard two-pass
% pdflatex compile.
%
% Usage in each Paper-NN.tex:
%   % =====================================================================
% TECT canonical LaTeX preamble (revtex4-2 / single-column / theorem-safe)
% =====================================================================
% Maintainer: Jusang Lee (jtkor@outlook.com)
% Binding from: 2026-05-06
% Purpose: a single source-of-truth preamble for every TECT paper
% (Paper-00 .. Paper-10 + Auxiliary notes). Designed to (a) eliminate
% font-corruption on pdflatex (T1 fontenc + lmodern), (b) make
% theorem-heavy mathematical-physics content render cleanly in a
% single-column layout (PRD class option, NOT PRL two-column), (c)
% provide consistent theorem numbering across all TECT papers via
% amsthm with a shared counter set, (d) make hyperref + cleveref
% cross-references resolve cleanly with the standard two-pass
% pdflatex compile.
%
% Usage in each Paper-NN.tex:
%   \input{../_shared/tect-preamble.tex}
%   \begin{document}
%   ...
%   \end{document}
%
% Build (every paper): `pdflatex Paper-NN && pdflatex Paper-NN`
% (the second pass resolves \ref{} / \cref{} forward references).
% A bash/PowerShell wrapper that automates this is at
% Docs/papers/papers/_shared/build-paper.sh / build-paper.ps1.
% =====================================================================

% ---- Document class --------------------------------------------------
% PRD class option of revtex4-2 with [reprint,onecolumn] gives the same
% APS heading / by-line / abstract style as PRL but in a wider single
% column that handles long display equations and theorem environments
% without column-break artefacts. nofootinbib keeps citations out of
% footnotes. The 11pt option restores standard font metrics; with T1
% fontenc + lmodern this gives non-bitmap glyphs.
\documentclass[%
  aps,prd,reprint,onecolumn,nofootinbib,11pt,%
  amsmath,amssymb,floatfix,longbibliography%
]{revtex4-2}

% ---- Encoding & fonts ------------------------------------------------
\usepackage[T1]{fontenc}        % proper 8-bit font encoding (no font corruption)
\usepackage[utf8]{inputenc}     % allow UTF-8 source
\usepackage{lmodern}            % scalable Latin Modern (replaces bitmap CM)
\usepackage{textcomp}           % extra text symbols
\usepackage{microtype}          % subtle kerning improvements

% ---- UTF-8 → LaTeX character mappings --------------------------------
% Maps the Unicode characters that occur in TECT body text (legacy from
% the chat-content auto-archival rule, CLAUDE.md §4) to LaTeX-safe
% commands. Without these declarations, T1+inputenc rejects the chars
% with "Unicode character X not set up for use with LaTeX".
% Adding declarations centrally (here) is preferable to per-file edits.
\DeclareUnicodeCharacter{00A7}{\S}              % §  (section sign)
\DeclareUnicodeCharacter{00B2}{\ensuremath{^{2}}}        % ²
\DeclareUnicodeCharacter{00B3}{\ensuremath{^{3}}}        % ³
\DeclareUnicodeCharacter{00B5}{\ensuremath{\mu}}         % µ (micro)
\DeclareUnicodeCharacter{00B7}{\ensuremath{\cdot}}       % · (middle dot)
\DeclareUnicodeCharacter{00D7}{\ensuremath{\times}}      % ×
\DeclareUnicodeCharacter{00E4}{\"a}             % ä
\DeclareUnicodeCharacter{00E9}{\'e}             % é
\DeclareUnicodeCharacter{00F6}{\"o}             % ö
\DeclareUnicodeCharacter{00FC}{\"u}             % ü
\DeclareUnicodeCharacter{010C}{\v{C}}           % Č
\DeclareUnicodeCharacter{010D}{\v{c}}           % č
\DeclareUnicodeCharacter{0394}{\ensuremath{\Delta}}      % Δ
\DeclareUnicodeCharacter{03A3}{\ensuremath{\Sigma}}      % Σ
\DeclareUnicodeCharacter{03A9}{\ensuremath{\Omega}}      % Ω
\DeclareUnicodeCharacter{03B1}{\ensuremath{\alpha}}      % α
\DeclareUnicodeCharacter{03B2}{\ensuremath{\beta}}       % β
\DeclareUnicodeCharacter{03B3}{\ensuremath{\gamma}}      % γ
\DeclareUnicodeCharacter{03B4}{\ensuremath{\delta}}      % δ
\DeclareUnicodeCharacter{03B5}{\ensuremath{\varepsilon}} % ε
\DeclareUnicodeCharacter{03BA}{\ensuremath{\kappa}}      % κ
\DeclareUnicodeCharacter{03BB}{\ensuremath{\lambda}}     % λ
\DeclareUnicodeCharacter{03BC}{\ensuremath{\mu}}         % μ
\DeclareUnicodeCharacter{03BD}{\ensuremath{\nu}}         % ν
\DeclareUnicodeCharacter{03C0}{\ensuremath{\pi}}         % π
\DeclareUnicodeCharacter{03C1}{\ensuremath{\rho}}        % ρ
\DeclareUnicodeCharacter{03C3}{\ensuremath{\sigma}}      % σ
\DeclareUnicodeCharacter{03C4}{\ensuremath{\tau}}        % τ
\DeclareUnicodeCharacter{03C6}{\ensuremath{\varphi}}     % φ
\DeclareUnicodeCharacter{03C7}{\ensuremath{\chi}}        % χ
\DeclareUnicodeCharacter{03C8}{\ensuremath{\psi}}        % ψ
\DeclareUnicodeCharacter{03C9}{\ensuremath{\omega}}      % ω
\DeclareUnicodeCharacter{2013}{--}              % – (en dash)
\DeclareUnicodeCharacter{2014}{---}             % — (em dash)
\DeclareUnicodeCharacter{2018}{`}               % ‘
\DeclareUnicodeCharacter{2019}{'}               % ’
\DeclareUnicodeCharacter{201C}{``}              % “
\DeclareUnicodeCharacter{201D}{''}              % ”
\DeclareUnicodeCharacter{2070}{\ensuremath{^{0}}}        % ⁰
\DeclareUnicodeCharacter{2071}{\ensuremath{^{i}}}        % ⁱ
\DeclareUnicodeCharacter{2122}{\textsuperscript{TM}}     % ™
\DeclareUnicodeCharacter{2126}{\ensuremath{\Omega}}      % Ω (ohm sign)
\DeclareUnicodeCharacter{210F}{\ensuremath{\hbar}}       % ℏ
\DeclareUnicodeCharacter{2115}{\ensuremath{\mathbb{N}}}  % ℕ
\DeclareUnicodeCharacter{2102}{\ensuremath{\mathbb{C}}}  % ℂ
\DeclareUnicodeCharacter{211D}{\ensuremath{\mathbb{R}}}  % ℝ
\DeclareUnicodeCharacter{2124}{\ensuremath{\mathbb{Z}}}  % ℤ
\DeclareUnicodeCharacter{211A}{\ensuremath{\mathbb{Q}}}  % ℚ
\DeclareUnicodeCharacter{2192}{\ensuremath{\to}}         % →
\DeclareUnicodeCharacter{21D2}{\ensuremath{\Rightarrow}} % ⇒
\DeclareUnicodeCharacter{2200}{\ensuremath{\forall}}     % ∀
\DeclareUnicodeCharacter{2203}{\ensuremath{\exists}}     % ∃
\DeclareUnicodeCharacter{2208}{\ensuremath{\in}}         % ∈
\DeclareUnicodeCharacter{2209}{\ensuremath{\notin}}      % ∉
\DeclareUnicodeCharacter{2227}{\ensuremath{\wedge}}      % ∧
\DeclareUnicodeCharacter{2228}{\ensuremath{\vee}}        % ∨
\DeclareUnicodeCharacter{2229}{\ensuremath{\cap}}        % ∩
\DeclareUnicodeCharacter{222A}{\ensuremath{\cup}}        % ∪
\DeclareUnicodeCharacter{2245}{\ensuremath{\cong}}       % ≅
\DeclareUnicodeCharacter{2248}{\ensuremath{\approx}}     % ≈
\DeclareUnicodeCharacter{2260}{\ensuremath{\neq}}        % ≠
\DeclareUnicodeCharacter{2264}{\ensuremath{\leq}}        % ≤
\DeclareUnicodeCharacter{2265}{\ensuremath{\geq}}        % ≥
\DeclareUnicodeCharacter{22C5}{\ensuremath{\cdot}}       % ⋅
\DeclareUnicodeCharacter{2295}{\ensuremath{\oplus}}      % ⊕
\DeclareUnicodeCharacter{2297}{\ensuremath{\otimes}}     % ⊗

% ---- Mathematics & symbols ------------------------------------------
\usepackage{amsmath,amssymb,bm,mathrsfs,mathtools}
\usepackage{amsthm}             % theorem environments (load BEFORE hyperref)

% ---- Theorem environments (shared counter set, TECT canonical) -------
% All theorems / lemmas / propositions / corollaries / remarks share a
% single counter so cross-references read consistently across a paper.
% Definitions get a separate counter (definitions are numbered in their
% own sequence by editorial convention).
\newtheorem{theorem}{Theorem}
\newtheorem{lemma}[theorem]{Lemma}
\newtheorem{proposition}[theorem]{Proposition}
\newtheorem{corollary}[theorem]{Corollary}

\theoremstyle{remark}
\newtheorem{remark}[theorem]{Remark}
\newtheorem{example}[theorem]{Example}

\theoremstyle{definition}
\newtheorem{definition}[theorem]{Definition}
\newtheorem{conjecture}[theorem]{Conjecture}
\newtheorem{hypothesis}[theorem]{Hypothesis}

% ---- Tables / floats / graphics --------------------------------------
\usepackage{booktabs}           % \toprule / \midrule / \bottomrule
\usepackage{array}
\usepackage{graphicx}
\usepackage{enumitem}           % \begin{itemize}[leftmargin=...] options
\usepackage{hyphenat}           % \hyp{} for breakable hyphens in \texttt
\usepackage{seqsplit}           % \seqsplit{} for long unbreakable strings
\sloppy                         % allow stretchier inter-word spacing to
                                % reduce overfull \hbox warnings on long URLs
                                % and long \texttt tokens (paragraph-level).
\emergencystretch=3em           % last-resort hbox stretch budget

% ---- Cross-references (hyperref last among the standard packages,
%      cleveref AFTER hyperref) ----------------------------------------
\usepackage[%
  colorlinks=true,%
  linkcolor=blue,%
  citecolor=blue,%
  urlcolor=blue,%
  breaklinks=true,%
  bookmarks=true,%
  bookmarksnumbered=true%
]{hyperref}
\usepackage[capitalise,nameinlink]{cleveref}

% ---- TECT-canonical author block macros -----------------------------
% (defined here so every paper has the same author / affiliation style)
\newcommand{\tectauthor}{%
  \author{Jusang Lee}%
  \email{jtkor@outlook.com}%
  \affiliation{Independent researcher, Republic of Korea}%
}

% ---- TECT-canonical archive footer ----------------------------------
\newcommand{\tectarchivefooter}{%
  \par\medskip
  \noindent\small\textit{Canonical archive}:
  \url{https://github.com/TECT-OpenPhysics/TECT}.\par%
}

% ---- TECT TODO / AUDIT-FLAG / HONEST-SCOPE inline marks --------------
% Used in drafts to highlight pending audit items without disturbing
% the body text style. Suppress via \tectsuppressmarks (define empty).
\providecommand{\tectauditflag}[1]{\textcolor{red}{\textbf{[AUDIT-FLAG: #1]}}}
\providecommand{\tecttodo}[1]{\textcolor{orange}{\textbf{[TODO: #1]}}}
\providecommand{\tecthonestscope}[1]{\textit{Honest scope: #1.}}

% =====================================================================
% End of TECT canonical preamble
% =====================================================================

%   \begin{document}
%   ...
%   \end{document}
%
% Build (every paper): `pdflatex Paper-NN && pdflatex Paper-NN`
% (the second pass resolves \ref{} / \cref{} forward references).
% A bash/PowerShell wrapper that automates this is at
% Docs/papers/papers/_shared/build-paper.sh / build-paper.ps1.
% =====================================================================

% ---- Document class --------------------------------------------------
% PRD class option of revtex4-2 with [reprint,onecolumn] gives the same
% APS heading / by-line / abstract style as PRL but in a wider single
% column that handles long display equations and theorem environments
% without column-break artefacts. nofootinbib keeps citations out of
% footnotes. The 11pt option restores standard font metrics; with T1
% fontenc + lmodern this gives non-bitmap glyphs.
\documentclass[%
  aps,prd,reprint,onecolumn,nofootinbib,11pt,%
  amsmath,amssymb,floatfix,longbibliography%
]{revtex4-2}

% ---- Encoding & fonts ------------------------------------------------
\usepackage[T1]{fontenc}        % proper 8-bit font encoding (no font corruption)
\usepackage[utf8]{inputenc}     % allow UTF-8 source
\usepackage{lmodern}            % scalable Latin Modern (replaces bitmap CM)
\usepackage{textcomp}           % extra text symbols
\usepackage{microtype}          % subtle kerning improvements

% ---- UTF-8 → LaTeX character mappings --------------------------------
% Maps the Unicode characters that occur in TECT body text (legacy from
% the chat-content auto-archival rule, CLAUDE.md §4) to LaTeX-safe
% commands. Without these declarations, T1+inputenc rejects the chars
% with "Unicode character X not set up for use with LaTeX".
% Adding declarations centrally (here) is preferable to per-file edits.
\DeclareUnicodeCharacter{00A7}{\S}              % §  (section sign)
\DeclareUnicodeCharacter{00B2}{\ensuremath{^{2}}}        % ²
\DeclareUnicodeCharacter{00B3}{\ensuremath{^{3}}}        % ³
\DeclareUnicodeCharacter{00B5}{\ensuremath{\mu}}         % µ (micro)
\DeclareUnicodeCharacter{00B7}{\ensuremath{\cdot}}       % · (middle dot)
\DeclareUnicodeCharacter{00D7}{\ensuremath{\times}}      % ×
\DeclareUnicodeCharacter{00E4}{\"a}             % ä
\DeclareUnicodeCharacter{00E9}{\'e}             % é
\DeclareUnicodeCharacter{00F6}{\"o}             % ö
\DeclareUnicodeCharacter{00FC}{\"u}             % ü
\DeclareUnicodeCharacter{010C}{\v{C}}           % Č
\DeclareUnicodeCharacter{010D}{\v{c}}           % č
\DeclareUnicodeCharacter{0394}{\ensuremath{\Delta}}      % Δ
\DeclareUnicodeCharacter{03A3}{\ensuremath{\Sigma}}      % Σ
\DeclareUnicodeCharacter{03A9}{\ensuremath{\Omega}}      % Ω
\DeclareUnicodeCharacter{03B1}{\ensuremath{\alpha}}      % α
\DeclareUnicodeCharacter{03B2}{\ensuremath{\beta}}       % β
\DeclareUnicodeCharacter{03B3}{\ensuremath{\gamma}}      % γ
\DeclareUnicodeCharacter{03B4}{\ensuremath{\delta}}      % δ
\DeclareUnicodeCharacter{03B5}{\ensuremath{\varepsilon}} % ε
\DeclareUnicodeCharacter{03BA}{\ensuremath{\kappa}}      % κ
\DeclareUnicodeCharacter{03BB}{\ensuremath{\lambda}}     % λ
\DeclareUnicodeCharacter{03BC}{\ensuremath{\mu}}         % μ
\DeclareUnicodeCharacter{03BD}{\ensuremath{\nu}}         % ν
\DeclareUnicodeCharacter{03C0}{\ensuremath{\pi}}         % π
\DeclareUnicodeCharacter{03C1}{\ensuremath{\rho}}        % ρ
\DeclareUnicodeCharacter{03C3}{\ensuremath{\sigma}}      % σ
\DeclareUnicodeCharacter{03C4}{\ensuremath{\tau}}        % τ
\DeclareUnicodeCharacter{03C6}{\ensuremath{\varphi}}     % φ
\DeclareUnicodeCharacter{03C7}{\ensuremath{\chi}}        % χ
\DeclareUnicodeCharacter{03C8}{\ensuremath{\psi}}        % ψ
\DeclareUnicodeCharacter{03C9}{\ensuremath{\omega}}      % ω
\DeclareUnicodeCharacter{2013}{--}              % – (en dash)
\DeclareUnicodeCharacter{2014}{---}             % — (em dash)
\DeclareUnicodeCharacter{2018}{`}               % ‘
\DeclareUnicodeCharacter{2019}{'}               % ’
\DeclareUnicodeCharacter{201C}{``}              % “
\DeclareUnicodeCharacter{201D}{''}              % ”
\DeclareUnicodeCharacter{2070}{\ensuremath{^{0}}}        % ⁰
\DeclareUnicodeCharacter{2071}{\ensuremath{^{i}}}        % ⁱ
\DeclareUnicodeCharacter{2122}{\textsuperscript{TM}}     % ™
\DeclareUnicodeCharacter{2126}{\ensuremath{\Omega}}      % Ω (ohm sign)
\DeclareUnicodeCharacter{210F}{\ensuremath{\hbar}}       % ℏ
\DeclareUnicodeCharacter{2115}{\ensuremath{\mathbb{N}}}  % ℕ
\DeclareUnicodeCharacter{2102}{\ensuremath{\mathbb{C}}}  % ℂ
\DeclareUnicodeCharacter{211D}{\ensuremath{\mathbb{R}}}  % ℝ
\DeclareUnicodeCharacter{2124}{\ensuremath{\mathbb{Z}}}  % ℤ
\DeclareUnicodeCharacter{211A}{\ensuremath{\mathbb{Q}}}  % ℚ
\DeclareUnicodeCharacter{2192}{\ensuremath{\to}}         % →
\DeclareUnicodeCharacter{21D2}{\ensuremath{\Rightarrow}} % ⇒
\DeclareUnicodeCharacter{2200}{\ensuremath{\forall}}     % ∀
\DeclareUnicodeCharacter{2203}{\ensuremath{\exists}}     % ∃
\DeclareUnicodeCharacter{2208}{\ensuremath{\in}}         % ∈
\DeclareUnicodeCharacter{2209}{\ensuremath{\notin}}      % ∉
\DeclareUnicodeCharacter{2227}{\ensuremath{\wedge}}      % ∧
\DeclareUnicodeCharacter{2228}{\ensuremath{\vee}}        % ∨
\DeclareUnicodeCharacter{2229}{\ensuremath{\cap}}        % ∩
\DeclareUnicodeCharacter{222A}{\ensuremath{\cup}}        % ∪
\DeclareUnicodeCharacter{2245}{\ensuremath{\cong}}       % ≅
\DeclareUnicodeCharacter{2248}{\ensuremath{\approx}}     % ≈
\DeclareUnicodeCharacter{2260}{\ensuremath{\neq}}        % ≠
\DeclareUnicodeCharacter{2264}{\ensuremath{\leq}}        % ≤
\DeclareUnicodeCharacter{2265}{\ensuremath{\geq}}        % ≥
\DeclareUnicodeCharacter{22C5}{\ensuremath{\cdot}}       % ⋅
\DeclareUnicodeCharacter{2295}{\ensuremath{\oplus}}      % ⊕
\DeclareUnicodeCharacter{2297}{\ensuremath{\otimes}}     % ⊗

% ---- Mathematics & symbols ------------------------------------------
\usepackage{amsmath,amssymb,bm,mathrsfs,mathtools}
\usepackage{amsthm}             % theorem environments (load BEFORE hyperref)

% ---- Theorem environments (shared counter set, TECT canonical) -------
% All theorems / lemmas / propositions / corollaries / remarks share a
% single counter so cross-references read consistently across a paper.
% Definitions get a separate counter (definitions are numbered in their
% own sequence by editorial convention).
\newtheorem{theorem}{Theorem}
\newtheorem{lemma}[theorem]{Lemma}
\newtheorem{proposition}[theorem]{Proposition}
\newtheorem{corollary}[theorem]{Corollary}

\theoremstyle{remark}
\newtheorem{remark}[theorem]{Remark}
\newtheorem{example}[theorem]{Example}

\theoremstyle{definition}
\newtheorem{definition}[theorem]{Definition}
\newtheorem{conjecture}[theorem]{Conjecture}
\newtheorem{hypothesis}[theorem]{Hypothesis}

% ---- Tables / floats / graphics --------------------------------------
\usepackage{booktabs}           % \toprule / \midrule / \bottomrule
\usepackage{array}
\usepackage{graphicx}
\usepackage{enumitem}           % \begin{itemize}[leftmargin=...] options
\usepackage{hyphenat}           % \hyp{} for breakable hyphens in \texttt
\usepackage{seqsplit}           % \seqsplit{} for long unbreakable strings
\sloppy                         % allow stretchier inter-word spacing to
                                % reduce overfull \hbox warnings on long URLs
                                % and long \texttt tokens (paragraph-level).
\emergencystretch=3em           % last-resort hbox stretch budget

% ---- Cross-references (hyperref last among the standard packages,
%      cleveref AFTER hyperref) ----------------------------------------
\usepackage[%
  colorlinks=true,%
  linkcolor=blue,%
  citecolor=blue,%
  urlcolor=blue,%
  breaklinks=true,%
  bookmarks=true,%
  bookmarksnumbered=true%
]{hyperref}
\usepackage[capitalise,nameinlink]{cleveref}

% ---- TECT-canonical author block macros -----------------------------
% (defined here so every paper has the same author / affiliation style)
\newcommand{\tectauthor}{%
  \author{Jusang Lee}%
  \email{jtkor@outlook.com}%
  \affiliation{Independent researcher, Republic of Korea}%
}

% ---- TECT-canonical archive footer ----------------------------------
\newcommand{\tectarchivefooter}{%
  \par\medskip
  \noindent\small\textit{Canonical archive}:
  \url{https://github.com/TECT-OpenPhysics/TECT}.\par%
}

% ---- TECT TODO / AUDIT-FLAG / HONEST-SCOPE inline marks --------------
% Used in drafts to highlight pending audit items without disturbing
% the body text style. Suppress via \tectsuppressmarks (define empty).
\providecommand{\tectauditflag}[1]{\textcolor{red}{\textbf{[AUDIT-FLAG: #1]}}}
\providecommand{\tecttodo}[1]{\textcolor{orange}{\textbf{[TODO: #1]}}}
\providecommand{\tecthonestscope}[1]{\textit{Honest scope: #1.}}

% =====================================================================
% End of TECT canonical preamble
% =====================================================================

%   \begin{document}
%   ...
%   \end{document}
%
% Build (every paper): `pdflatex Paper-NN && pdflatex Paper-NN`
% (the second pass resolves \ref{} / \cref{} forward references).
% A bash/PowerShell wrapper that automates this is at
% Docs/papers/papers/_shared/build-paper.sh / build-paper.ps1.
% =====================================================================

% ---- Document class --------------------------------------------------
% PRD class option of revtex4-2 with [reprint,onecolumn] gives the same
% APS heading / by-line / abstract style as PRL but in a wider single
% column that handles long display equations and theorem environments
% without column-break artefacts. nofootinbib keeps citations out of
% footnotes. The 11pt option restores standard font metrics; with T1
% fontenc + lmodern this gives non-bitmap glyphs.
\documentclass[%
  aps,prd,reprint,onecolumn,nofootinbib,11pt,%
  amsmath,amssymb,floatfix,longbibliography%
]{revtex4-2}

% ---- Encoding & fonts ------------------------------------------------
\usepackage[T1]{fontenc}        % proper 8-bit font encoding (no font corruption)
\usepackage[utf8]{inputenc}     % allow UTF-8 source
\usepackage{lmodern}            % scalable Latin Modern (replaces bitmap CM)
\usepackage{textcomp}           % extra text symbols
\usepackage{microtype}          % subtle kerning improvements

% ---- UTF-8 → LaTeX character mappings --------------------------------
% Maps the Unicode characters that occur in TECT body text (legacy from
% the chat-content auto-archival rule, CLAUDE.md §4) to LaTeX-safe
% commands. Without these declarations, T1+inputenc rejects the chars
% with "Unicode character X not set up for use with LaTeX".
% Adding declarations centrally (here) is preferable to per-file edits.
\DeclareUnicodeCharacter{00A7}{\S}              % §  (section sign)
\DeclareUnicodeCharacter{00B2}{\ensuremath{^{2}}}        % ²
\DeclareUnicodeCharacter{00B3}{\ensuremath{^{3}}}        % ³
\DeclareUnicodeCharacter{00B5}{\ensuremath{\mu}}         % µ (micro)
\DeclareUnicodeCharacter{00B7}{\ensuremath{\cdot}}       % · (middle dot)
\DeclareUnicodeCharacter{00D7}{\ensuremath{\times}}      % ×
\DeclareUnicodeCharacter{00E4}{\"a}             % ä
\DeclareUnicodeCharacter{00E9}{\'e}             % é
\DeclareUnicodeCharacter{00F6}{\"o}             % ö
\DeclareUnicodeCharacter{00FC}{\"u}             % ü
\DeclareUnicodeCharacter{010C}{\v{C}}           % Č
\DeclareUnicodeCharacter{010D}{\v{c}}           % č
\DeclareUnicodeCharacter{0394}{\ensuremath{\Delta}}      % Δ
\DeclareUnicodeCharacter{03A3}{\ensuremath{\Sigma}}      % Σ
\DeclareUnicodeCharacter{03A9}{\ensuremath{\Omega}}      % Ω
\DeclareUnicodeCharacter{03B1}{\ensuremath{\alpha}}      % α
\DeclareUnicodeCharacter{03B2}{\ensuremath{\beta}}       % β
\DeclareUnicodeCharacter{03B3}{\ensuremath{\gamma}}      % γ
\DeclareUnicodeCharacter{03B4}{\ensuremath{\delta}}      % δ
\DeclareUnicodeCharacter{03B5}{\ensuremath{\varepsilon}} % ε
\DeclareUnicodeCharacter{03BA}{\ensuremath{\kappa}}      % κ
\DeclareUnicodeCharacter{03BB}{\ensuremath{\lambda}}     % λ
\DeclareUnicodeCharacter{03BC}{\ensuremath{\mu}}         % μ
\DeclareUnicodeCharacter{03BD}{\ensuremath{\nu}}         % ν
\DeclareUnicodeCharacter{03C0}{\ensuremath{\pi}}         % π
\DeclareUnicodeCharacter{03C1}{\ensuremath{\rho}}        % ρ
\DeclareUnicodeCharacter{03C3}{\ensuremath{\sigma}}      % σ
\DeclareUnicodeCharacter{03C4}{\ensuremath{\tau}}        % τ
\DeclareUnicodeCharacter{03C6}{\ensuremath{\varphi}}     % φ
\DeclareUnicodeCharacter{03C7}{\ensuremath{\chi}}        % χ
\DeclareUnicodeCharacter{03C8}{\ensuremath{\psi}}        % ψ
\DeclareUnicodeCharacter{03C9}{\ensuremath{\omega}}      % ω
\DeclareUnicodeCharacter{2013}{--}              % – (en dash)
\DeclareUnicodeCharacter{2014}{---}             % — (em dash)
\DeclareUnicodeCharacter{2018}{`}               % ‘
\DeclareUnicodeCharacter{2019}{'}               % ’
\DeclareUnicodeCharacter{201C}{``}              % “
\DeclareUnicodeCharacter{201D}{''}              % ”
\DeclareUnicodeCharacter{2070}{\ensuremath{^{0}}}        % ⁰
\DeclareUnicodeCharacter{2071}{\ensuremath{^{i}}}        % ⁱ
\DeclareUnicodeCharacter{2122}{\textsuperscript{TM}}     % ™
\DeclareUnicodeCharacter{2126}{\ensuremath{\Omega}}      % Ω (ohm sign)
\DeclareUnicodeCharacter{210F}{\ensuremath{\hbar}}       % ℏ
\DeclareUnicodeCharacter{2115}{\ensuremath{\mathbb{N}}}  % ℕ
\DeclareUnicodeCharacter{2102}{\ensuremath{\mathbb{C}}}  % ℂ
\DeclareUnicodeCharacter{211D}{\ensuremath{\mathbb{R}}}  % ℝ
\DeclareUnicodeCharacter{2124}{\ensuremath{\mathbb{Z}}}  % ℤ
\DeclareUnicodeCharacter{211A}{\ensuremath{\mathbb{Q}}}  % ℚ
\DeclareUnicodeCharacter{2192}{\ensuremath{\to}}         % →
\DeclareUnicodeCharacter{21D2}{\ensuremath{\Rightarrow}} % ⇒
\DeclareUnicodeCharacter{2200}{\ensuremath{\forall}}     % ∀
\DeclareUnicodeCharacter{2203}{\ensuremath{\exists}}     % ∃
\DeclareUnicodeCharacter{2208}{\ensuremath{\in}}         % ∈
\DeclareUnicodeCharacter{2209}{\ensuremath{\notin}}      % ∉
\DeclareUnicodeCharacter{2227}{\ensuremath{\wedge}}      % ∧
\DeclareUnicodeCharacter{2228}{\ensuremath{\vee}}        % ∨
\DeclareUnicodeCharacter{2229}{\ensuremath{\cap}}        % ∩
\DeclareUnicodeCharacter{222A}{\ensuremath{\cup}}        % ∪
\DeclareUnicodeCharacter{2245}{\ensuremath{\cong}}       % ≅
\DeclareUnicodeCharacter{2248}{\ensuremath{\approx}}     % ≈
\DeclareUnicodeCharacter{2260}{\ensuremath{\neq}}        % ≠
\DeclareUnicodeCharacter{2264}{\ensuremath{\leq}}        % ≤
\DeclareUnicodeCharacter{2265}{\ensuremath{\geq}}        % ≥
\DeclareUnicodeCharacter{22C5}{\ensuremath{\cdot}}       % ⋅
\DeclareUnicodeCharacter{2295}{\ensuremath{\oplus}}      % ⊕
\DeclareUnicodeCharacter{2297}{\ensuremath{\otimes}}     % ⊗

% ---- Mathematics & symbols ------------------------------------------
\usepackage{amsmath,amssymb,bm,mathrsfs,mathtools}
\usepackage{amsthm}             % theorem environments (load BEFORE hyperref)

% ---- Theorem environments (shared counter set, TECT canonical) -------
% All theorems / lemmas / propositions / corollaries / remarks share a
% single counter so cross-references read consistently across a paper.
% Definitions get a separate counter (definitions are numbered in their
% own sequence by editorial convention).
\newtheorem{theorem}{Theorem}
\newtheorem{lemma}[theorem]{Lemma}
\newtheorem{proposition}[theorem]{Proposition}
\newtheorem{corollary}[theorem]{Corollary}

\theoremstyle{remark}
\newtheorem{remark}[theorem]{Remark}
\newtheorem{example}[theorem]{Example}

\theoremstyle{definition}
\newtheorem{definition}[theorem]{Definition}
\newtheorem{conjecture}[theorem]{Conjecture}
\newtheorem{hypothesis}[theorem]{Hypothesis}

% ---- Tables / floats / graphics --------------------------------------
\usepackage{booktabs}           % \toprule / \midrule / \bottomrule
\usepackage{array}
\usepackage{graphicx}
\usepackage{enumitem}           % \begin{itemize}[leftmargin=...] options
\usepackage{hyphenat}           % \hyp{} for breakable hyphens in \texttt
\usepackage{seqsplit}           % \seqsplit{} for long unbreakable strings
\sloppy                         % allow stretchier inter-word spacing to
                                % reduce overfull \hbox warnings on long URLs
                                % and long \texttt tokens (paragraph-level).
\emergencystretch=3em           % last-resort hbox stretch budget

% ---- Cross-references (hyperref last among the standard packages,
%      cleveref AFTER hyperref) ----------------------------------------
\usepackage[%
  colorlinks=true,%
  linkcolor=blue,%
  citecolor=blue,%
  urlcolor=blue,%
  breaklinks=true,%
  bookmarks=true,%
  bookmarksnumbered=true%
]{hyperref}
\usepackage[capitalise,nameinlink]{cleveref}

% ---- TECT-canonical author block macros -----------------------------
% (defined here so every paper has the same author / affiliation style)
\newcommand{\tectauthor}{%
  \author{Jusang Lee}%
  \email{jtkor@outlook.com}%
  \affiliation{Independent researcher, Republic of Korea}%
}

% ---- TECT-canonical archive footer ----------------------------------
\newcommand{\tectarchivefooter}{%
  \par\medskip
  \noindent\small\textit{Canonical archive}:
  \url{https://github.com/TECT-OpenPhysics/TECT}.\par%
}

% ---- TECT TODO / AUDIT-FLAG / HONEST-SCOPE inline marks --------------
% Used in drafts to highlight pending audit items without disturbing
% the body text style. Suppress via \tectsuppressmarks (define empty).
\providecommand{\tectauditflag}[1]{\textcolor{red}{\textbf{[AUDIT-FLAG: #1]}}}
\providecommand{\tecttodo}[1]{\textcolor{orange}{\textbf{[TODO: #1]}}}
\providecommand{\tecthonestscope}[1]{\textit{Honest scope: #1.}}

% =====================================================================
% End of TECT canonical preamble
% =====================================================================


\begin{document}

\title{A Pillar~6 electroweak-bridge programme note for TECT: first
$\Delta F<0$ broken-energy data point at $N=16$ within the
continuum-limit-scan protocol}

\author{Jusang Lee}
\email{jtkor@outlook.com}
\affiliation{Independent researcher, Republic of Korea}

\date{\today}

\begin{abstract}
We present a candidate electroweak-bridge programme exploring the
embedding of the BCC shell structure of the topological condensate
into the Standard-Model three-generation $+$ Higgs-doublet sector
within Topological Energy Condensate Theory (TECT). Using the
continuum-limit-scan protocol (Math235), we evaluate the candidate
electroweak-bridge projector $\mathcal{P}_{\mathrm{EW}}$ from the
$12$-dimensional BCC shell moduli space $V_{\mathrm{shell}}$ into the
candidate $2$-dimensional Higgs-doublet support $V_{\mathrm{EW}}$
(Math234). At lattice resolutions $N \in \{4,8,16,32\}$, the
broken-phase free-energy gap $\Delta F = F_{\mathrm{broken}}-
F_{\mathrm{symmetric}}$ first turns negative at $N=16$ with
$\Delta F_{N=16} \approx -324.94$ in TECT internal units; the $N=32$
value $\Delta F_{N=32}\approx -1021.7$ extends the trend but is
\emph{not} a continuum endpoint. The current canonical Pillar~6 tier
is \textbf{T4 STRONG EVIDENCE}. Under the Math292 four-gauge
acceptance criterion (G1 finite-$f$; G2 broken-energy
$\Delta F<0$; G3 transverse-Hessian PSD;
G4 extraction \texttt{OK}), the present data point fulfils G2 only
at $N=16$. \emph{G3 verdict at $N=32$, $\mu^2=-0.7$ is
\textbf{INDETERMINATE} under the current $3$-translation
zero-mode projector} (Math315 first transverse-Lanczos
diagnostic, 2026-05-06): $\mathrm{ARPACK}$ smallest-algebraic
eigenvalue solver returns $0/1$ converged after
$\sim 10^{4}$ matvecs, the canonical near-zero clustered-spectrum
saddle-manifold signature predicted in Math290. G3 (and G1, G4)
remain pending the Math315 \S 4 follow-up sequence
(shift-invert / LOBPCG / deeper Newton convergence) before any
re-certification of Math310-AddA. Math293 false-negative taxonomy and the
Math294/Math294-AddA basin-of-attraction analyses motivate $N=16$ as
the empirical first-entry resolution but do not, at this tier, certify
embedding. A pre-registered F-Pillar6 falsification gate (deadline
\textbf{2026-05-29}) governs the next $\mathrm{T4}\to\mathrm{T6}$
promotion attempt; the present document does \emph{not} certify
Pillar~6 closure.
\end{abstract}

\maketitle

% =====================================================================
\section{Introduction}\label{sec:intro}
% =====================================================================
The Standard Model's three-fermion-generation $+$ Higgs-doublet
structure is, within most frameworks, an external phenomenological
input. TECT poses a different question: can the structure be made
\emph{compatible} with a topological-condensate vacuum that supports
the requisite gauge embedding? The present paper does not assert that
the answer is positive at theorem level. It reports the first
numerical data point at which the candidate broken phase becomes
energetically favoured ($\Delta F<0$) within the continuum-limit-scan
protocol, and audits the data point against the Math292 four-gauge
acceptance criterion.

The BCC condensate is not uniform; its $12$ nearest-neighbour
generators form a crystallographic shell. Each shell mode is a
candidate degree of freedom for the embedding programme; the
projection of the $12$-dimensional shell moduli onto a candidate
$2$-dimensional Higgs-doublet support is the object of study (Math234).
The present paper does \emph{not} prove that the three SM generations
correspond to specific shell modes, nor that the Yukawa hierarchy is
fixed; both remain Pillar~6 \textbf{T1 OPEN} (Math234 sub-tasks
SO(10)-rep and Yukawa-extraction).

The contribution of this paper is purely the $\Delta F$ data point and
its honest audit against Math292 acceptance: the broken phase exists
at $N\geq 16$, the embedding remains pending.

% =====================================================================
\section{Setting and main result}\label{sec:setup}
% =====================================================================
The BCC shell is described by the moduli vector
$\phi = (\phi_1,\dots,\phi_{12})\in V_{\mathrm{shell}}$, where each
$\phi_i\in\mathbb{C}$ is the amplitude on one BCC nearest-neighbour
direction. The order parameter is
\begin{equation}\label{eq:OP}
  \Psi(\mathbf{r}) \;=\; \sum_{i=1}^{12} \phi_i \exp(i\mathbf{q}_i\cdot\mathbf{r}),
\end{equation}
with $\mathbf{q}_i$ the $12$ BCC nearest-neighbour wave vectors.

The candidate electroweak-bridge projector
$\mathcal{P}_{\mathrm{EW}}: V_{\mathrm{shell}}\to V_{\mathrm{EW}}\cong
\mathbb{C}^2$ is defined by coefficient decomposition (Math234):
\begin{equation}\label{eq:PEW}
  H \;=\; \mathcal{P}_{\mathrm{EW}}\cdot\phi \;=\;
  \begin{pmatrix} \alpha_H \\ \beta_H \end{pmatrix},
  \qquad
  \alpha_H = \sum_{j\in S_u} c_j^\alpha \phi_j,
  \quad
  \beta_H = \sum_{j\in S_d} c_j^\beta \phi_j,
\end{equation}
where $S_u,S_d$ partition the $12$ directions, and
$c_j^{\alpha/\beta}$ are structure constants fixed by the BCC lattice
geometry. The map \eqref{eq:PEW} is a candidate projector; whether it
identifies the physical Higgs doublet of the SM is the embedding
question that remains open.

\begin{theorem}[First $\Delta F<0$ data point at $N=16$
(\textbf{T4 STRONG EVIDENCE}, G2-of-4-gauges only)]
\label{thm:main}
Under hypotheses
\begin{itemize}
  \item[(\textbf{H1})] the continuum-limit-scan protocol of
    Math235--236;
  \item[(\textbf{H2})] the Math292 four-gauge acceptance criterion
    \emph{interpreted in the present paper at the G2 (broken-energy
    entry) gauge only}, i.e.\ $\Delta F<0$, with G1 (finite $f$),
    G3 (transverse-Hessian PSD), G4 (extraction \texttt{OK}) marked
    pending per Math310-AddA;
  \item[(\textbf{H3})] the Math293 false-negative taxonomy
    classifying $N<16$ data as sub-critical;
  \item[(\textbf{H4})] the Math294/Math294-AddA basin-of-attraction
    analyses identifying $N=16$ as the empirical marginal-basin
    resolution,
\end{itemize}
the broken-phase free-energy gap evaluated at $N=16$ is
\begin{equation}\label{eq:deltaF}
  \boxed{\;
    \Delta F_{N=16} \;\approx\; -324.94 \pm 8.7
    \quad\text{(TECT units; G2 satisfied; G1+G4 pending; G3 \emph{INDETERMINATE} per Math315).}
  \;}
\end{equation}
\end{theorem}

\begin{remark}[Honest scope]\label{rmk:honest-scope}
\eqref{eq:deltaF} is a single-resolution, single-gauge data point
within the Math292 four-gauge framework. \emph{It does not certify
embedding.} Promotion to T6 \textsc{proved conditional} requires
all four Math292 gauges to pass simultaneously at $N \geq 16$, plus
discharge of the Math310-AddA transverse-projection re-certification.
The F-Pillar6 falsification gate (deadline 2026-05-29) governs the
next $\mathrm{T4}\to\mathrm{T6}$ promotion attempt.
\end{remark}

% =====================================================================
\section{Method: continuum-limit scan and the Math292 four-gauge audit}
\label{sec:method}
% =====================================================================
\paragraph{Step 1 (continuum-limit-scan protocol).}
Math235 defines a four-stage protocol implemented in Math236
\cite{Math236}: (i) discretise the BCC shell on lattices with
$N\in\{4,8,16,32\}$; (ii) minimise the free-energy functional
$F[\phi]$ numerically on each grid; (iii) compute
$\Delta F = F_{\mathrm{broken}}-F_{\mathrm{symmetric}}$ as a function
of $N$; (iv) extrapolate to the continuum ($N\to\infty$) via a
quadratic fit, conditional on at least three super-critical data
points and a verified residual.

\paragraph{Step 2 (Math292 four-gauge acceptance criterion).}
The Math292 acceptance gate for an embedding-data point is the
\emph{conjunction} of four conditions:
\begin{equation}\label{eq:gauges}
  \mathrm{Math292\!-\!Accept} \;\;\Longleftrightarrow\;\;
  \underbrace{(f > 0)}_{\textbf{G1}} \;\wedge\;
  \underbrace{(\Delta F < 0)}_{\textbf{G2}} \;\wedge\;
  \underbrace{(\lambda_{\min}^{\mathrm{transverse}} \geq 0)}_{\textbf{G3}}
  \;\wedge\;
  \underbrace{(\text{extraction} = \texttt{OK})}_{\textbf{G4}}.
\end{equation}
The label \emph{A-valid} that appeared in earlier drafts of this paper
referred only to G2; that usage was an over-claim and is corrected
here. The present data point passes G2 at $N=16$ only.

\paragraph{G3 status under the current projector (Math315
INDETERMINATE).} The first explicit transverse-Lanczos diagnostic at
$N=32$, $\mu^2=-0.7$ on the
\texttt{math236\_A1p0\_seeded\_N32\_resumed/Psi\_best\_F.npy}
candidate (Math315, 2026-05-06) recovered \emph{zero} converged
$\lambda_{\min}^{\mathrm{transverse}}$ eigenvalues after
$\sim 10^{4}$ ARPACK matvecs at \texttt{which='SA'}, with the
asymmetry-probe ratio $|(J - J^T)v|/(|Jv| + |J^Tv|) \sim 10^{-5}$
ruling out an operator implementation defect. This is the canonical
densely-clustered-near-zero spectrum signature predicted by Math290
\S Step~5: the broken-phase configuration sits on a saddle manifold
of dimension $\geq 7$ (3 translations + 1 global $U(1)$ phase + at
least $3$ internal Cartan / shell phases), of which the current
\texttt{build\_zero\_mode\_projector} default projects out only the
$3$ translations. The G3 verdict is therefore recorded as
\textbf{INDETERMINATE}, \emph{not} FAIL, pending the Math315 \S 4
follow-up programme: (N1) shift-invert near zero with phase-mode
projection and larger Krylov subspace, (N2) LOBPCG, or
(N3) deeper Newton convergence to
$\|R\|/\sqrt{\mathrm{dof}} < 10^{-5}$. G1 (finite $f$) and G4
(extraction \texttt{OK}) remain pending the same Math310-AddA
re-certification programme.

\paragraph{Step 3 (false-negative taxonomy).}
Math293 classifies sub-critical data ($\Delta F\geq 0$ at $N<16$ due
to finite-size artefact) as false negatives that must not block the
search. Conversely, $\Delta F<0$ at $N\geq 16$ counts as a candidate
positive provided the remaining Math292 gauges are independently
verified \cite{Math293}.

\paragraph{Step 4 (basin-of-attraction analyses).}
Math294 and Math294-AddA analyse the marginal basin of the broken
phase and identify $N=16$ as the empirical first-entry resolution.
This justifies $N=16$ as the operating point of the present audit; it
does not promote the data point past T4 in the absence of the G1/G3/G4
verifications.

\paragraph{Step 5 (numerical audit).}
Math290 documents wrapper and solver bugs encountered in the first
production run, with corrections applied to produce the values
reported here \cite{Math290}.

% =====================================================================
\section{Numerical results and Math292 gauge audit}\label{sec:results}
% =====================================================================
The continuum-limit-scan results for $\Delta F$ in TECT internal
units are summarised in Table~\ref{tab:dfscan}.

\begin{table}[h]
\centering
\begin{tabular}{cccc}
\toprule
$N$ & $\Delta F$ (numerical) & Phase & Math292 verdict \\
\midrule
$4$  & $+127.3$  & symmetric & G2 fails (sub-critical) \\
$8$  & $+45.2$   & symmetric & G2 fails (sub-critical) \\
$16$ & $-324.94$ & broken    & G2 only; G1/G3/G4 pending \\
$32$ & $-1021.7$ & broken    & G2 only; G1/G3/G4 pending \\
\bottomrule
\end{tabular}
\caption{Continuum-limit-scan data for the broken-phase free-energy
gap. Only G2 (broken-energy entry) of the Math292 four-gauge
acceptance criterion is satisfied at $N=16,32$; the remaining gauges
G1, G3, G4 await Math310-AddA re-certification.}
\label{tab:dfscan}
\end{table}

\begin{remark}[$N=32$ is not a continuum endpoint]\label{rmk:N32}
The $N=32$ value $\Delta F \approx -1021.7$ extends the broken-phase
trend but is \emph{not} a continuum-limit verdict. A continuum-limit
extrapolation requires (i) a third super-critical data point at
$N\geq 64$, (ii) a quadratic-fit residual below the per-point
numerical-error budget, and (iii) the Math292 gauges G1/G3/G4 verified
at every super-critical $N$ in the fit range. None of (i)--(iii) is
discharged by the present data set; the $N=32$ value is therefore
reported as a plateau-trending data point and \emph{not} as a
confirmed endpoint.
\end{remark}

The candidate Higgs-doublet support $(\alpha_H,\beta_H)$ extracted at
$N=16$ inherits its symbolic structure from Math234. Whether this
extraction reproduces the physical SM Higgs charge and hypercharge
assignments is a separate Math234-extraction sub-task and is, at the
present canonical tier, \textbf{T1 OPEN}.

% =====================================================================
\section{Discussion and next steps}\label{sec:discussion}
% =====================================================================
The present paper records a single empirical fact: within the
continuum-limit-scan protocol of Math235--236, the broken-phase
free-energy gap $\Delta F$ first turns negative at lattice resolution
$N=16$, with values $\Delta F_{N=16}\approx -324.94$ and
$\Delta F_{N=32}\approx -1021.7$ in TECT internal units. The
data point passes the G2 (broken-energy entry) gauge of the Math292
four-gauge acceptance criterion only, and remains pending on
G1/G3/G4. It does not, at the present canonical tier, certify
the candidate electroweak-bridge embedding programme.

Residual tasks before $\mathrm{T4}\to\mathrm{T6}$ promotion:
\begin{enumerate}
  \item \textbf{Math310-AddA transverse-projection re-certification}
    discharging G3 ($\lambda_{\min}^{\mathrm{transverse}}\geq 0$);
  \item independent verification of G1 (finite $f$) and G4
    (extraction \texttt{OK}) at $N\in\{16,32\}$;
  \item a third super-critical data point at $N\geq 64$ with
    Math292 four-gauge passes, enabling the continuum-limit
    quadratic fit;
  \item separate Math234-extraction sub-task verifying the candidate
    Higgs-doublet identification.
\end{enumerate}
The pre-registered F-Pillar6 falsification gate carries deadline
\textbf{2026-05-29}. Success on all four tasks supports the
$\mathrm{T4}\to\mathrm{T6}$ promotion attempt; partial success keeps
the result at T4 with explicit caveat enumeration.

% =====================================================================
\begin{acknowledgments}
The author thanks the numerical team for the continuum-limit-scan
runs, and Claude Opus 4.7 for audit-taxonomy classification (Math293)
and basin-of-attraction confirmation (Math294, Math294-AddA). Wrapper
errors in the first production run were caught and corrected per
Math290. This work is part of TECT, canonical archive at
\url{https://github.com/TECT-OpenPhysics/TECT}.
\end{acknowledgments}

% =====================================================================
\bibliographystyle{apsrev4-2}
\bibliography{../_shared/tect-references}

\end{document}
